\begin{center}
{\heitifont\zihao{2} 摘\quad 要}
\end{center}

视觉语言模型在退化图像场景下的语义理解能力不足是制约其实际部署的关键瓶颈之一。脑电图作为一种高时间分辨率的神经成像手段，记录了被试观看视觉刺激时的大脑皮层活动，其中可能包含可被深度学习模型利用的视觉语义信息。本课程设计系统研究了在小样本 EEG-image 配对数据条件下，真实 EEG 信号能否在视觉退化时为视觉语言模型提供有效的语义增强。

本设计以 CLIP 预训练语义空间为跨模态对齐基础，将图像嵌入、类别文本原型和 EEG 表示统一到同一向量空间中。在此基础上，设计并实现了 A2 时间-频谱-空间三支路 EEG 编码器，通过并行提取 EEG 信号的时域动态、频域分布和通道空间特征，经由门控残差融合机制对退化图像的视觉嵌入进行动态补充。进一步提出 VTF token 级融合方法，利用视觉-EEG 交叉注意力在 CLIP ViT visual token 层面注入脑电信息，生成增强视觉表示。最终将增强视觉 token 接入 Qwen2-VL-2B 生成式视觉语言模型，采用 LoRA 低秩适配（$r=8$）进行高效微调，并结合多候选重排序策略输出自由形式图像描述。

实验基于 Thought2Text 数据集，采用 image-level split 防止信息泄漏，设置六种视觉退化条件和四种 EEG 对照模式。在三种子 × 六种条件的 18 组评估中，A2 语义融合模型的真实 EEG 均优于 vision-only、shuffled EEG 和 random EEG 对照组。在生成式任务中，真实 EEG 的 caption 类别命中率在强退化条件下较纯视觉提升约 8--9 个百分点，多候选重排序下的有效输出率达到 97.8\%。实验结果初步表明，真实配对 EEG 在视觉退化条件下能够为视觉语言模型提供可验证的语义辅助信息。

\vspace{0.5cm}
\noindent\textbf{关键词：}脑电信号；视觉语言模型；CLIP；视觉退化；多模态融合；LoRA；图像描述生成

\clearpage
\begin{center}
{\bfseries\zihao{2} Abstract}
\end{center}

This course design investigates EEG-assisted semantic recognition and caption generation under degraded visual conditions. A CLIP-based semantic bridge is established, followed by an A2 temporal-spectral-spatial EEG encoder with gated residual fusion. A VTF token-level cross-attention module injects EEG information into CLIP ViT visual tokens. The enhanced tokens are then processed by Qwen2-VL with LoRA adaptation and best-of-N reranking. Systematic experiments across six degradation types with real/shuffled/random EEG controls demonstrate consistent improvements from paired EEG signals.

\vspace{0.5cm}
\noindent\textbf{Keywords:} EEG; Vision-Language Model; CLIP; Visual Degradation; Multimodal Fusion; LoRA; Image Captioning
